\begin{solapartat}
\punts{0,2} $\vec{F}_m = q\left(\vec{v}\times\vec{B}\right);\ F_m = qvB\sin(90^\circ) = qvB$

\punts{0,2} $\left\{\begin{aligned} F_m &= m a_c\\ qvB &= m\,\frac{v^2}{R} \end{aligned}\right.$

\punts{0,3}
\figura[0.2]{sol_fig1.png}

\punts{0,2} $\omega = \dfrac{7\times 2\pi}{1,29\times 10^{-3}} = 3,41\times 10^{4}\ \mathrm{rad\,s^{-1}}$

\punts{0,2} $qvB = m\,\dfrac{v^2}{r} \Rightarrow qB = m\,\dfrac{v}{r} = m\omega \Rightarrow m = \dfrac{qB}{\omega}$

% DUBTE: les puntuacions parcials de la pauta sumen 1,3 p (0,2 + 0,2 + 0,3 + 0,2 + 0,2 + 0,2), però l'apartat val 1 punt.
\punts{0,2} $m = \dfrac{\left|-1,60\times 10^{-19}\right|\times 45,0\times 10^{-3}}{3,41\times 10^{4}} = 2,11\times 10^{-25}\ \mathrm{kg}$
\end{solapartat}

\begin{solapartat}
\punts{0,1} $\vec{F}_m = q\left(\vec{v}\times\vec{B}\right)$

\punts{0,2} $F_m = qvB\sin(90^\circ) = qvB$

\punts{0,4} $B = \dfrac{8,00\times 10^{-14}}{1,60\times 10^{-19}\times 5,00\times 10^{5}} = 1,00\ \mathrm{T}$

\punts{0,3} Direcció perpendicular al pla del paper, sentit cap dins del pla del paper. La justificació es pot fer a partir de les propietats del producte vectorial o aplicant la regla de la ma dreta.
\figura[0.2]{sol_fig2.png}
\end{solapartat}
